\documentclass[11pt]{article}
\usepackage[T1]{fontenc}
\usepackage{textcomp}
\usepackage[margin=0.75in]{geometry}
\usepackage{amsmath,amssymb,amsthm}
\usepackage{array,booktabs}
\usepackage{hyperref}
\setlength{\parindent}{0pt}
\newtheorem{theorem}{Theorem}
\theoremstyle{definition}
\newtheorem{definition}{Definition}
\title{Flow Proof Artifact: Comb-choose-zero-n}
\author{Generated by \texttt{flow doc proof}}
\date{Flow Proof Book}
\begin{document}
\maketitle
Choosing from zero items gives zero when k is positive.\\[0.5em]
\textbf{Source.} graham-knuth-patashnik — *Concrete Mathematics*\\[0.5em]
\bigskip
\noindent{\large \textbf{Derived fact 1.} \textit{choose(0, k) = 0 for k > 0} \hfill Comb \cdot choose \cdot choosing from empty gives zero when k is positive}\par
\smallskip\noindent\textit{Coordinate: Comb \cdot choose \cdot choosing from empty gives zero when k is positive}\par
\smallskip\noindent\textit{Source: graham-knuth-patashnik}\par
\smallskip\noindent\textit{Built on: choosing none gives one, for choose on Comb, pascal recurrence holds, for choose on Comb}\par
\medskip
\noindent\textbf{Goal.} choose(0, k) = 0 for k > 0\par
\noindent$\displaystyle \forall k \in \mathbb{N}\quad choose(0, k) = 0$\par
\medskip
\noindent
\begin{tabular}{@{} >{\raggedright\arraybackslash}p{0.06\textwidth} >{\raggedright\arraybackslash}p{0.47\textwidth} >{\raggedleft\arraybackslash}p{0.41\textwidth} @{}}
\textbf{\#} & \textbf{Proof} & \textbf{Mathematics} \\
\midrule
\textbf{1.} & We prove that choosing from empty gives zero when k is positive for choose on Comb. & \\[0.45em]
\textbf{2.} & We split into exhaustive cases — the claim must hold in each one. & \\[0.45em]
\textbf{3.} & Case 1 (see step 2): suppose k is zero. & \\[0.45em]
\textbf{4.} & We invoke the definitional clause governing choose on Comb: choosing none gives one, for choose on Comb (instantiated for 0). & \\[0.45em]
\textbf{5.} & From step 3 and step 4, this implies choose(0, k) equals 0 in this case. & $\displaystyle choose(0, k) = 0$ \\[0.45em]
\textbf{6.} & Case 2 (see step 2): neither disjunct holds. & \\[0.45em]
\textbf{7.} & Let j denote the predecessor of k. & \\[0.45em]
\textbf{8.} & We invoke the definitional clause governing choose on Comb: pascal recurrence holds, for choose on Comb (instantiated for 0, k). & \\[0.45em]
\textbf{9.} & From step 6, step 7, and step 8, this implies choose(0, k) equals 0. Together with the other cases (step 3 and step 6), the goal is discharged. Hence proven. & $\displaystyle choose(0, k) = 0$ \\[0.45em]
\bottomrule
\end{tabular}
\label{thm:Comb-choose-choosing_from_empty_gives_zero_when_k_is_positive}
\par\medskip\hrule
\end{document}
